% ============================================================
%  Corriges TD Chapitre 1 : Les ensembles
% ============================================================
\chapter{Les ensembles}
\chaptermark{Les ensembles}

\begin{corrige}{1.1}
\begin{enumerate}
\item En extension : $A = \{2, 3, 4, 5, 6, 7, 8, 9, 10\}$, donc $|A| = 9$ éléments.
\item Le mot « licence » compte sept lettres, mais la lettre c et la lettre e s'y répètent ; les lettres distinctes sont $l$, $i$, $c$, $e$ et $n$, donc $B = \{l, i, c, e, n\}$ et $|B| = 5$.
\item L'entier $4$ est bien compris entre $2$ et $10$, donc $4 \in A$ est \emph{vrai} ; l'entier $11$ dépasse $10$, donc $11 \notin A$ est \emph{vrai} ; l'entier $3$ est compris entre $2$ et $10$, donc $3 \in A$ est \emph{vrai}.
\item Pour un produit cartésien, $|A \times B| = |A| \times |B| = 9 \times 5 = 45$ couples. Un couple s'écrit $(x, y)$ avec $x \in A$ et $y \in B$ ; par exemple $(7, i)$, ou encore $(2, n)$.
\end{enumerate}
\end{corrige}

\begin{corrige}{1.2}
\begin{enumerate}
\item L'union réunit les éléments des deux ensembles : $A \cup B = \{1, 2, 3, 4, 5, 6\}$. L'intersection ne garde que les éléments communs : $A \cap B = \{3, 4\}$. Le complémentaire de $A$ dans $U$ regroupe les éléments absents de $A$ : $\overline{A} = \{5, 6, 7, 8\}$ ; de même $\overline{B} = \{1, 2, 7, 8\}$.
\item $\overline{A} \cap \overline{B} = \{7, 8\}$ : seuls $7$ et $8$ sont absents de $A$ \emph{et} de $B$. L'union $A \cup B = \{1, 2, 3, 4, 5, 6\}$ couvre six éléments, donc son complémentaire vaut $\overline{A \cup B} = \{7, 8\}$. Les deux résultats coïncident : c'est la première loi de De Morgan, $\overline{A \cup B} = \overline{A} \cap \overline{B}$.
\item $\overline{A \cap B} = U \setminus \{3, 4\} = \{1, 2, 5, 6, 7, 8\}$, et $\overline{A} \cup \overline{B} = \{5, 6, 7, 8\} \cup \{1, 2, 7, 8\} = \{1, 2, 5, 6, 7, 8\}$. Les deux résultats coïncident encore : c'est la seconde loi de De Morgan, $\overline{A \cap B} = \overline{A} \cup \overline{B}$. Nier une union revient à nier chaque membre puis intersecter ; nier une intersection revient à nier chaque membre puis unir.
\end{enumerate}
\end{corrige}

\begin{corrige}{1.3}
\begin{enumerate}
\item On commence par la zone centrale : $|A \cap B| = 15$. La zone « bureautique seulement » vaut $48 - 15 = 33$, et la zone « algorithmique seulement » vaut $39 - 15 = 24$. Le diagramme porte donc $33$, $15$ et $24$ dans les trois zones des deux cercles.
\item Les étudiants suivant au moins un module occupent les trois zones des cercles : $33 + 15 + 24 = 72$ étudiants.
\item Il reste $90 - 72 = 18$ étudiants hors des deux cercles, dans la zone $\overline{A \cup B}$ : aucun des deux modules.
\item La formule donne $|A \cup B| = |A| + |B| - |A \cap B| = 48 + 39 - 15 = 72$, et l'on retrouve bien les $72$ étudiants de la question 2 : les $15$ étudiants des deux modules, comptés deux fois dans $48 + 39$, ont été retirés une fois.
\end{enumerate}
\end{corrige}

\begin{corrige}{1.4}
\begin{enumerate}
\item On remplit du centre vers l'extérieur. Noyau : $|P \cap J \cap S| = 4$. Intersections doubles seulement : $P \cap J$ sans SQL vaut $12 - 4 = 8$ ; $P \cap S$ sans JavaScript vaut $9 - 4 = 5$ ; $J \cap S$ sans Python vaut $7 - 4 = 3$. Zones simples : Python seul vaut $34 - 8 - 5 - 4 = 17$ ; JavaScript seul vaut $27 - 8 - 3 - 4 = 12$ ; SQL seul vaut $22 - 5 - 3 - 4 = 10$.
\item L'union des trois ensembles couvre $17 + 12 + 10 + 8 + 5 + 3 + 4 = 59$ développeurs.
\item Il en reste $70 - 59 = 11$ qui ne pratiquent aucun des trois langages.
\item Somme des huit régions : $17 + 12 + 10 + 8 + 5 + 3 + 4 + 11 = 70$, ce qui est conforme. La méthode du centre vers l'extérieur évite de compter plusieurs fois les développeurs qui pratiquent plusieurs langages.
\end{enumerate}
\end{corrige}

\begin{corrige}{1.5}
\textbf{Expression $E$.} La distributivité de $\cup$ sur $\cap$ met $B$ en facteur :
\[
E = (A \cap B) \cup (\overline{A} \cap B) = (A \cup \overline{A}) \cap B = U \cap B = B ,
\]
avec la complémentarité $A \cup \overline{A} = U$ puis le neutre $U \cap B = B$. \emph{Contrôle.} Avec $U = \{1, 2, 3, 4\}$, $A = \{1, 2\}$ et $B = \{2, 3\}$ : $A \cap B = \{2\}$, $\overline{A} = \{3, 4\}$ donc $\overline{A} \cap B = \{3\}$, d'où $E = \{2, 3\} = B$. Vérifié.

\textbf{Expression $F$.} Cette fois la distributivité de $\cup$ sur $\cap$ factorise par $A$ :
\[
F = (A \cup B) \cap (A \cup \overline{B}) = A \cup (B \cap \overline{B}) = A \cup \varnothing = A ,
\]
avec $B \cap \overline{B} = \varnothing$ puis le neutre $A \cup \varnothing = A$. \emph{Contrôle.} Avec les mêmes ensembles : $A \cup B = \{1, 2, 3\}$, $\overline{B} = \{1, 4\}$ donc $A \cup \overline{B} = \{1, 2, 4\}$, et l'intersection des deux vaut $\{1, 2\} = A$. Vérifié.

Dans les deux cas, l'expression de départ, qui mêle trois opérations, se réduit à un seul ensemble : la lecture par le diagramme de Venn l'annonçait, la distributivité l'écrit, et le test numérique le confirme.
\end{corrige}
